% !TEX root = ../analisi-dei-capitolati.tex
% Capitolo 5 - C0 - AIRA

\section{C0 -- AIRA}\label{sec:c0}

\begin{scheda}
  \voce{Proponente}{Sanmarco Informatica S.p.A.}
  \voce{Dominio}{automazione delle richieste B2B multicanale}
\end{scheda}

\subsection{Obiettivo}
Le richieste dei clienti B2B (ordini, ticket, appuntamenti, documenti
contabili) arrivano da canali eterogenei: e-mail con allegati, note
vocali, messaggi, foto di cataloghi. Il progetto chiede uno strumento che
monitori i canali, interpreti le richieste con trascrizione vocale,
\glo{OCR} e un LLM, ne estragga dati strutturati e li inoltri in
automatico ai gestionali aziendali tramite API REST o e-mail.

\subsection{Requisiti principali}
Sono obbligatori il controllo continuo dei canali, un filtro di sicurezza
in ingresso, l'elaborazione semantica tramite LLM, l'avvio di azioni
configurabili e un'interfaccia amministrativa web; sono opzionali una
dashboard di monitoraggio e le risposte automatiche al mittente. Le
tecnologie consigliate sono Angular, Java con Spring Boot, LangChain4j,
Whisper o Tesseract per voce e immagini e un database relazionale.

\subsection{Incontro con il proponente}
All'incontro del 21 ottobre 2026 (VE~2026-10-21) Sanmarco Informatica ha
chiarito che non ci saranno costi esterni (modelli \emph{open weight} in
locale, per esempio Qwen), che WhatsApp e i gestionali saranno simulati e
che i canali prioritari sono e-mail e audio. Ha inoltre chiesto misure di
sicurezza oltre la whitelist: antispam, offuscamento e cifratura dei dati,
isolamento dei tenant. Il supporto prevede un incontro ogni due settimane,
poi settimanale dopo il PoC.

\subsection{Valutazione del gruppo}

\begin{fortedeboli}
  \fd{\textbf{Problema concreto} e diffuso nelle aziende.}
     {\textbf{Ampiezza tecnologica}: voce, OCR, LLM, integrazioni e
      interfaccia web insieme.}
  \fd{\textbf{Nessun costo}: modelli locali e canali simulati.}
     {\textbf{Limiti dei modelli locali}, legati all'hardware dei membri.}
  \fd{\textbf{Priorità chiarite} e supporto regolare.}
     {\textbf{Sicurezza estesa} su un nucleo già ampio.}
\end{fortedeboli}

C0 è la seconda alternativa del gruppo: concreto e senza costi, ma il più
ampio dei tre finalisti sul piano tecnologico e il meno votato.
